\begin{solapartat}
\punts{0,5} Hem d'imposar la conservació del nombre de nucleons i de la càrrega elèctrica. Igualment, com que s'emet un leptó (electró) que és partícula, s'ha d'emetre un leptó que sigui antipartícula, per tant s'emetrà un antineutrí:
\[ {}^{14}_{\phantom{1}6}C \to {}^{14}_{\phantom{1}7}N + {}^{\phantom{-}0}_{-1}e + {}^{0}_{0}\bar{\nu}_e \]

\punts{0,25} Es tracta d'una desintegració beta ($\beta^-$) perquè s'emet un electró.

\punts{0,5} L'equació de desintegració és $N(t) = N_0\,e^{-\lambda t}$ en què $\lambda$ és la constant de desintegració. El període de semidesintegració $t_{1/2}$ és aquell en el qual la mostra radioactiva $N$ s'ha reduït a la meitat. Per a aquest temps es compleix:

$N(t) = \frac{N_0}{2} = N_0\,e^{-\lambda t_{1/2}}$ i per tant, $\frac{1}{2} = e^{-\lambda t_{1/2}}$ i obtenim $\lambda\cdot t_{1/2} = -ln\frac{1}{2}$. Aïllant el període de semidesintegració tenim: $t_{1/2} = \frac{ln2}{\lambda}$.
\end{solapartat}

\begin{solapartat}
\punts{0,25} Utilitzem l'expressió $t_{1/2} = \frac{ln2}{\lambda}$ per calcular la constant de desintegració:
\[ \lambda = \frac{ln2}{t_{1/2}} = \frac{ln2}{5730} = 1,21\cdot 10^{-4}\,anys^{-1} \]
o el seu equivalent en segons $\lambda = 3,83\cdot 10^{-12}\,s^{-1}$

\punts{0,5} Per calcular quin percentatge de carboni-14 roman a la mostra respecte el que hi havia en el moment de la mort de l'animal, tornem a utilitzar l'equació de desintegració $N(t) = N_0\,e^{-\lambda t}$. Aplicant l'expressió pel temps que ha passat obtenim el quocient entre el carboni-14 actual i l'inicial (moment de la mort de l'animal):
\[ \frac{N}{N_0} = e^{-\lambda t} = e^{-1,21\cdot 10^{-4}\cdot 6010} = 0,483\,. \]
Per tant, queda un 48,3\% del carboni-14 inicial.

\punts{0,5} La variació de l'activitat amb el temps ens la dona la relació següent: $A(t) = A_0\,e^{-\lambda t}$ Si l'activitat actual és de 3,63 Bq i han passat 6010 anys, l'activitat de l'os en el moment de la mort de l'animal, $A_0$, és: $A_0 = \frac{A(t)}{e^{-\lambda t}} = \frac{3,63}{0,483} = 7,51\ \mathrm{Bq}$.
\end{solapartat}
